\chapter{Introduction}

\section{Présentation personnelle}
Je suis Sam FOKAM WABO, stagiaire en formation au titre professionnel d'enseignant de la conduite et de la sécurité routière (ECSR). Mon parcours professionnel antérieur, notamment dans le développement informatique, m'a appris à structurer un projet, à rechercher des informations, à analyser des besoins et à documenter les décisions prises. Dans le cadre de ma reconversion vers l'enseignement de la conduite, je souhaite mobiliser cette méthode de travail au service d'une pédagogie active et de la prévention des risques routiers.

Le certificat de compétences professionnelles n\textsuperscript{o} 2 (CCP2) porte sur la sensibilisation de l'ensemble des usagers de la route à l'adoption de comportements sûrs et respectueux de l'environnement. Il ne s'agit pas seulement de transmettre des règles : l'enjeu est de comprendre un public, de construire une intervention pertinente, de l'animer puis d'en évaluer les effets et sa propre pratique.

Cette action est conduite en binôme avec Manane SADIK, également stagiaire en formation ECSR. Notre commanditaire est le centre AFPA de Nice, représenté dans nos échanges par Mme Muriel LIÉGEOIS, responsable de formation du centre AFPA Nice/Cannes. L'intervention, prévue le 4 décembre 2026, s'adresse à un groupe de stagiaires en formation de plomberie. Nous avons retenu le thème général « Sur la route du chantier », centré sur l'alcool et les distracteurs pendant les déplacements professionnels.

\section{Démarche et périmètre du présent mémoire}
Ce mémoire retrace une démarche progressive : recherche d'un commanditaire, clarification de sa demande, enquête auprès du public, hiérarchisation des besoins, conception d'une réponse et recherche de ressources pédagogiques. Il est rédigé à la première personne pour rendre explicites mes initiatives, mes choix et les analyses effectuées conjointement avec ma collègue.

\textbf{État du dossier au 11 novembre 2026 :} la phase de diagnostic et une première version de la préparation pédagogique sont documentées. L'action du 4 décembre n'a pas encore eu lieu. Ses résultats, les retours des participants et l'analyse réflexive ne sont donc pas anticipés : ils seront rédigés après l'intervention, sur la base d'observations et de pièces vérifiables.
